\documentclass[10pt,a4paper]{amsart}
\usepackage[margin=19mm]{geometry}
\usepackage{amsmath,amssymb,mathtools}
\usepackage[scr=rsfs,cal=euler]{mathalfa}
\usepackage{xfrac}
\usepackage[unicode,hidelinks,bookmarks=false]{hyperref}
\newcommand{\ie}{i.e.\ }
\newcommand{\cf}{cf.\ }
\newcommand{\resp}{respectively\ }
\newcommand{\Subsection}{\subsection}
\providecommand{\nobreakdash}{\nobreak\hbox{-}}
\setlength{\emergencystretch}{2em}
\multlinegap0pt
\usepackage{bm}
\usepackage{bbm}
\usepackage{dsfont}
\usepackage{xspace}
\usepackage{cancel}
\usepackage{mathtools}
\usepackage{amsthm}
\makeatletter
\@ifundefined{theorem}{\newtheorem{theorem}{Theorem}}{}
\makeatother
\title[Critical lack of equilibrium in stochastic kinetic proofread]{Critical lack of equilibrium in stochastic kinetic proofreading}
\author{E. Franco, J. J. L. Vel\'azquez}
\date{Source: arXiv:2508.19983v1 (CC BY 4.0)}
\begin{document}
\maketitle

\noindent\emph{Source and licence.} Critical lack of equilibrium in stochastic kinetic proofreading. Version arXiv:2508.19983v1 (CC BY 4.0), distributed under the Creative Commons Attribution 4.0 International licence, as recorded in the arXiv licence metadata for this exact version and shown on the abstract page https://arxiv.org/abs/2508.19983v1. Full rights notice: licenses/arxiv-2508.19983-CC-BY-4.0.txt.\par\smallskip

\noindent\textit{Editorial scope.} Counted unit: the Theorem whose source label is \texttt{theo:half line asymptotics delta supercrit} in the article (source0.tex lines 1492--1505, source label \texttt{theo:half line asymptotics delta supercrit}), delivered together with the article's own complete original proof of it (source0.tex lines 1506--1639, one \texttt{proof} environment of 134 lines). No mathematics is written, repaired or re-derived: statement and proof are the article's own text line for line. Required context retained uncounted: theo:half line asymptotics delta subcrit (lines 1320--1327); a (lines 1383--1385); Phi (lines 1415--1417). Every definition, display and label of those lines is retained unchanged. Omitted from this package and not redistributed: 1--1319, 1328--1382, 1386--1414, 1418--1491, 1640--2234. Nothing inside a retained excerpt is omitted or altered: every retained body is delivered exactly as the article's lines are written, so no omission receipt of any kind accompanies this package. Citations: the retained excerpts carry no citation command at all, so no bibliography entry of the article is transcribed and no reference is redistributed. Reference adaptation: none was needed and none was made; every reference inside the delivered unit resolves inside it, so no unresolved reference or citation is produced. Printed slips kept literal: no printed word, symbol, hypothesis or number of the counted statement or of its proof was altered. Layout and class: the article's own class is \texttt{article} with packages including geometry, inputenc, hyperref, caption, booktabs, graphicx, latexsym, lipsum, authblk, babel, color, dsfont, float, bbm; the wrapper is the lane's pinned \texttt{amsart} editorial wrapper at 10pt on A4, which restates the theorem-like environments the retained text needs and adds the article's own macro definitions that the retained text needs (none), with extra wrapper packages (bm, bbm, dsfont, xspace, cancel, mathtools). The delivered numbering is the unit's own. Assets: no retained span contains a figure environment, a graphics-inclusion command, a PDF-page inclusion, a table or a TikZ picture; no image file is copied into this package, the package declares no asset dependency and no image file is redistributed with it. Licence: version arXiv:2508.19983v1 of this article is distributed under the Creative Commons Attribution 4.0 International licence, as recorded in the arXiv licence metadata for this exact version and shown on the abstract page https://arxiv.org/abs/2508.19983v1. A full rights notice is delivered at licenses/arxiv-2508.19983-CC-BY-4.0.txt. Characters: every retained excerpt is pure ASCII, so the delivered engine, XeLaTeX with its default Latin Modern text font, reports no missing character.
\begin{theorem} \label{theo:half line asymptotics delta subcrit}
Let us define $\tau:= 2 \alpha e^{\frac{E+ \Delta }{2}} t$ and let $k:=\lfloor \theta \tau \rfloor$ for  $\theta \geq 0$.
\\  Assume that $\Delta < \Delta_c $. 
Then 
    \begin{equation} \label{asympt half subc}
    \lim_{t \to \infty } \frac{ n_{k} (t) }{ e^{-k E} } = A(0) \overline n_S.
    \end{equation}
\end{theorem}

\begin{equation} \label{a}
A(z)=a(w):= \frac{(w-w_A)^{-1}}{ 2 \alpha e^{(E+\Delta)/2 } }  \ \text{ and }  \ \hat{n}_S(z)= \tilde{n}_S (w): = \frac{\left(w- \cosh\left( \frac{E- \Delta }{2} \right) \right) (w-w_0)^{-1} (w-w_S)^{-1}}{2\alpha e^{(E+\Delta )/2 }} . 
\end{equation}

with \begin{equation} \label{Phi}
\Phi(w):= w - \theta  \log(w + \sqrt{w^2-1}). 
\end{equation}

\begin{theorem} \label{theo:half line asymptotics delta supercrit}
Let us define $\tau:= 2 \alpha e^{\frac{E+ \Delta }{2}} t$ and let $k:=\lfloor \theta \tau \rfloor$ for  $\theta \geq 0$.
    Assume that $\Delta > \Delta_c $. 
    \begin{enumerate} 
\item If $\theta $ is such that $w_\theta < w_S $, then we have that  
     \begin{equation}\label{asympt half superc thetha small}
    \lim_{t \to \infty } \frac{ n_{k} (t) }{ e^{k( \log(\varphi(\sigma))- E ) } } = A(0) B(0) \overline n_S  .
    \end{equation}
 \item    If instead $\theta $ is such that  $w_\theta  >w_S  $, then 
     \begin{equation}\label{asympt half superc thetha large}
    \lim_{t \to \infty } \frac{ n_{k} (t) }{ e^{k( \log(\varphi(\sqrt{1+\theta^2 }))- E ) } } = 0.
    \end{equation}
    \end{enumerate}
\end{theorem}

\begin{proof}
As in the proof of Theorem \ref{theo:half line asymptotics delta subcrit} we apply the inverse Laplace transform formula to deduce that
\[
n_k(t)= \frac{1}{2 \pi i } e^{- k E } \int_{\delta - i \infty }^{\delta + i \infty } e^{ zt } A(z) \hat{n}_S(z) \left( 1+ B(z)  \varphi(z)^k  \right) dz   
\]
where $z_A > \delta > 0 $. 
As in the proof of Theorem \ref{theo:half line asymptotics delta subcrit} $n_k=\ell_k+ h_k$ where  
\begin{equation} \label{lk}
\ell_k (t):= \frac{1}{2 \pi i } e^{- k E } \int_{\delta - i \infty }^{\delta + i \infty } e^{ zt } A(z) \hat{n}_S(z)  dz = e^{- k E } \left( A(0) \overline {n}_S + e^{ z_S t } A(z_S ) \frac{1}{z_S} \right)
\end{equation}
and 
\begin{equation}\label{def:h_k2}
h_k (t):= \frac{1}{2 \pi i } e^{- k E } \int_{\delta - i \infty }^{\delta + i \infty } e^{ zt } A(z) \hat{n}_S(z) B(z) \varphi(z)^k   dz
\end{equation}
Notice that since $ z_S <0 $ the dominant term in $ \ell_k(t)$ as $t \rightarrow \infty $ is $ e^{- k E } A(0) \overline {n}_S $. Hence
\[
\ell_{\tau \theta } \left( \frac{\tau }{ 2 \alpha e^{(E+\Delta)/2} }\right)  \sim  e^{- \tau \theta E - \tau w_0 } e^{ \tau w_0 } A(0) \overline {n}_S  \text{ as } \tau \to \infty. 
\]
Moreover 
\begin{align*}
h_k(t)&=   \frac{C(\alpha, E, \Delta)}{2 \pi i } e^{k \frac{\Delta - E }{2}}e^{- w_0 \tau }g_k(\tau ) 
\end{align*}
where $g_k$ and $C(\alpha, E, \Delta)$  are defined as in the proof of Theorem \ref{theo:half line asymptotics delta subcrit}. 
With the same arguments as the ones used in the proof of Theorem \ref{theo:half line asymptotics delta subcrit} we deduce that 
\[
g_{\tau \theta } (\tau)  \sim  \int_{\gamma}f(w) dw -  \frac{    \sqrt{2\pi} i \theta a(w_\theta ) \tilde{n}_S(w_\theta )}{(\theta+ \sqrt{1+ \theta^2 } - e^{(E-\Delta)/2}) (1+\theta^2)^{1/4}} \frac{ e^{ \tau \Phi(w_\theta ) }}{\sqrt{\tau}} \text{ as } \tau \to \infty. 
\]
where $\gamma  $ is as in the proof of Theorem \ref{theo:half line asymptotics delta subcrit}. 
In order to compute $\int_{\gamma}f(w) dw$ we need to distinguish between the following three cases. 
\begin{enumerate} 
\item $w_0 > w_\theta > w_S  $, 
\item $ w_0 > w_S > w_\theta $, 
\item $ w_\theta > w_0> w_S $.
\end{enumerate}


\paragraph{Case 1: $w_0 > w_\theta > w_S  $.}
Assume that 1. holds, hence that $\theta $ is such that $w_S < w_\theta < w_0 $.  
Then we have that 
\begin{align*}
 \int_{\gamma}f(w) dw=  2 \pi i  \operatorname{Res}_{w=w_0} f(w).
\end{align*} 
Now notice that
\[
\operatorname{Res}_{w=w_0} f(w) = \operatorname{Res}_{w=w_0} \tilde{n}_S (w)  a(w_0) b(w_0) e^{\tau \Phi (w_0) } 
\]
where we are using the notation 
\[ 
b(w):= \frac{1}{w+\sqrt{w^2-1}-e^{(E-\Delta)/2}}
\] 
and where the function $a$ is given by \eqref{a} and the function $\Phi$ by \eqref{Phi}. 
Therefore 

\[
g_{\tau \theta } (\tau)  \sim \operatorname{Res}_{w=w_0} \tilde{n}_S (w)  a(w_0) b(w_0) e^{\tau \Phi (w_0) }  -  \frac{    \sqrt{2\pi} i \theta a(w_\theta ) \tilde{n}_S(w_\theta )}{(\theta+ \sqrt{1+ \theta^2 } - e^{(E-\Delta)/2}) (1+\theta^2)^{1/4}} \frac{ e^{ \tau \Phi(w_\theta ) }}{\sqrt{\tau}} \text{ as } \tau \to \infty. 
\]
Hence
\begin{align*}
        h_{\theta \tau } (t) \sim &C(\alpha, E, \Delta)e^{ - \theta \tau E - \tau w_0 } [   b(w_0)   a(w_0) e^{ \frac{ \theta \tau (\Delta +E) }{2}}  e^{\tau   \Phi(w_0)} \operatorname{Res}_{w=w_0} \tilde{n}_S (w) \\
        &- \frac{   \theta  b(w_\theta)   a(w_\theta ) \tilde{n}_S(w_\theta )  
        e^{ \tau \Phi(w_\theta ) + \tau\theta  \frac{E+ \Delta}{2}}
        }{\sqrt{2 \pi \tau } (1+\theta^2)^{1/4}} ]   \text{ as } \tau \to \infty. 
\end{align*}

In order to study the behaviour of  $n_{\theta \tau}(t)$ as $\tau \to \infty $ we need to identify the leading order term in $\ell_{\theta \tau}+ h_{\theta \tau}$.
Since $\Delta > \Delta_c $ and $w_0 > w_\theta  $, we have that
\[  
\Phi(w_0) + \theta\frac{E+ \Delta }{2}=\max\left\{ w_0 , \Phi(w_0) + \theta\frac{E+ \Delta }{2}, \Phi(w_\theta) + \theta \frac{E+ \Delta }{2}  \right\}. 
\] 
Indeed, the fact that $\Delta > \Delta_c $ implies that 
\[
\Phi(w_0) + \theta\frac{E+ \Delta }{2} >w_0 > w_S. 
\]
Moreover notice that by the definition of $w_\theta $ we have that  $\Phi(w_0) > \Phi(w_\theta)$. Hence the desired conclusion follows. We conclude that in this case 
\begin{equation} \label{asympt n 1b}
    n_{\theta \tau} (t) \sim C(\alpha, \Delta, E)   b(w_0) a(w_0) \operatorname{Res}_{w=w_0} \tilde{n}_S (w)e^{ \frac{ \theta \tau (\Delta + E) }{2}}  e^{\tau   \Phi(w_0)}  \text{ as } \tau \to \infty. 
\end{equation}
Therefore using that 
\[ 
C(\alpha, \Delta, E)   b(w_0) a(w_0) \operatorname{Res}_{w=w_0} \tilde{n}_S (w)e^{ \frac{ \theta \tau (\Delta + E) }{2}}  e^{\tau   \Phi(w_0)} = -  A(0) B(0) \overline{n}_S
\]
we deduce that \eqref{asympt half superc thetha small} holds. 

\paragraph{Case 2: $w_0 > w_S > w_\theta   $. }
The main difference between case $1$ and case $2$ is the value of the integral $\int_\gamma f(w( dw ) $. 
Indeed, when  $\theta $ is such that $w_S > w_\theta $ we have that 
\begin{align*}
 \int_{\gamma} f(w) dw=  2 \pi i \left( \operatorname{Res}_{w=w_0} f(w)+ \operatorname{Res}_{w=w_S} f(w) \right) . 
\end{align*}
As a consequence, computations similar to the ones made for case 1 imply that 
\begin{align*}
        h_{\theta \tau } (t)\sim &e^{ - \theta \tau E - \tau w_0 } [   b(w_0)  \operatorname{Res}_{w=w_0} \tilde{n}_S (w) e^{ \frac{ \theta \tau (\Delta - E) }{2}}  e^{\tau   \Phi(w_0)} +b(w_S) \operatorname{Res}_{n=n_S} \left(\tilde{n}_S (w ) \right) e^{ \frac{ \theta \tau (\Delta - E) }{2}}  e^{\tau   \Phi(w_S)} \\
        &
        - \frac{  \theta  b(w_\theta)   a(w_\theta ) \tilde{n}_S(w_\theta )}{\sqrt{2 \pi } (1+\theta^2)^{1/4}} \frac{ e^{ \tau \Phi(w_\theta ) + \tau\theta  \frac{E+ \Delta}{2}}}{\sqrt{\tau}}].
\end{align*}



In order to find the leading order term in $\ell_{\theta \tau } + h_{\theta \tau }$
 we prove that
\begin{equation}\label{comparison 2.b}
\Phi(w_0) + \theta\frac{E+ \Delta }{2}=\max\left\{ w_0 , \Phi(w_0) + \theta\frac{E+ \Delta }{2}, \Phi(w_\theta) + \theta \frac{E+ \Delta }{2}, \Phi(w_S) + \theta\frac{E+ \Delta }{2}  \right\}.
\end{equation}
Indeed, the fact that $\Delta > \Delta_c$ implies that 
\[
\Phi(w_0) +\theta\frac{E+ \Delta }{2} >w_0 > w_S > \Phi(w_s) \text{ and }  \Phi(w_0)> \Phi(w_\theta).
\]
Therefore the asymptotics \eqref{asympt half superc thetha small} follow. 

\paragraph{Case 3: $w_\theta  > w_0 > w_S   $. }
When  $\theta $ is such that $ w_\theta  > w_0 > w_S  $ we have that 
\begin{align*}
 \int_{\gamma} f(w) dw= 0. 
\end{align*}
As a consequence, 
\begin{align*}
        h_{\theta \tau } (t)\sim 
        - \frac{ C(\alpha, E, \Delta)  \theta  b(w_\theta)   a(w_\theta ) \tilde{n}_S(w_\theta )}{\sqrt{2 \pi } (1+\theta^2)^{1/4}} \frac{ e^{ \tau \Phi(w_\theta ) + \tau\theta  \frac{E+ \Delta}{2}}}{\sqrt{\tau}} \text{ as } \tau \to \infty. 
\end{align*}

Notice that the fact that $ w_\theta > w_0 > w_S  $ and $\Delta > \Delta_c$ imply that
\[ 
\Phi(w_\theta) + \theta \frac{E+ \Delta }{2} =\max\left\{ w_0 , \Phi(w_\theta) + \theta \frac{E+ \Delta }{2} \right\}. 
\]
Hence
\[ 
   n_{\theta \tau } (t)\sim 
        - \frac{ C(\alpha, E, \Delta)  \theta  b(w_\theta)   a(w_\theta ) \tilde{n}_S(w_\theta )}{\sqrt{2 \pi } (1+\theta^2)^{1/4}} \frac{ e^{ \tau \Phi(w_\theta ) + \tau\theta  \frac{E+ \Delta}{2}}}{\sqrt{\tau}} \text{ as } \tau \to \infty
        \text{ as } \tau \to \infty
\]
Therefore 
\eqref{asympt half superc thetha large} follows. 

\end{proof}

\end{document}
